% Vietnamese translation for OI Public Library
% Source-PDF: USACO/【由 BirDOR 整理】USACO 2001-2007 月赛测试数据+题目+题解/USACO月赛题解.pdf
\documentclass[11pt]{article}
\usepackage{oipl-vn}
\usepackage{array}
\usepackage{booktabs}
\usepackage{longtable}

\title{Lời giải các kỳ USACO Monthly Contests 2002--2008}
\author{}
\date{}

\newcolumntype{L}[1]{>{\raggedright\arraybackslash}p{#1}}
\newenvironment{solutiontable}{%
  \small
  \setlength{\tabcolsep}{4pt}
  \renewcommand{\arraystretch}{1.14}
  \begin{longtable}{L{0.19\linewidth}L{0.39\linewidth}L{0.34\linewidth}}
  \toprule
  \textbf{Tên bài} & \textbf{Mô tả ngắn} & \textbf{Thuật toán / ghi chú}\\
  \midrule
  \endfirsthead
  \toprule
  \textbf{Tên bài} & \textbf{Mô tả ngắn} & \textbf{Thuật toán / ghi chú}\\
  \midrule
  \endhead
}{%
  \bottomrule
  \end{longtable}
  \normalsize
}
\newcommand{\problemrow}[3]{\textbf{#1} & #2 & #3\\\midrule}

\begin{document}
\maketitle

\origtitle{USACO 2002--2008 Monthly Contest Solution Notes}
\sourcepath{USACO/monthly-contest-solutions.pdf}

\section{USACO 2002 February}

\begin{solutiontable}
\problemrow{Fiber Communications}{Có \(N\) người đứng thành vòng tròn và \(M\) cặp người muốn liên lạc. Mỗi lần chỉ được nối hai người kề nhau. Hỏi cần nối ít nhất bao nhiêu cạnh để thỏa mọi cặp liên lạc.}{Liệt kê vị trí cắt vòng, rồi làm tương tự tô màu trên đường thẳng; dùng DSU để gộp các đoạn cần nối.}
\problemrow{Power Hungry Cows}{Ban đầu có \(X\) và \(1\). Mỗi bước chọn hai số hiện có, có thể trùng nhau, rồi dùng nhân hoặc chia để thay một trong hai số. Hỏi ít nhất bao nhiêu bước để tạo \(X^P\).}{Bài toán tương đương tạo số mũ \(P\) bằng cộng/trừ các số mũ đã có. Có thể BFS, chỉ mở rộng khi số lớn không vượt \(50000\) và số nhỏ không vượt \(100\); cũng có thể dùng IDA*.}
\problemrow{Cow Cycling}{Có \(N\) con bò chạy \(D\) vòng; ban đầu mỗi con có thể lực \(E\). Nếu chạy \(X\) vòng, bò dẫn đầu mất \(X^2\) thể lực mỗi phút, các bò sau mất \(X\). Cần thời gian ít nhất để chạy xong.}{DP: \(F[i][j][k]\) là thời gian ít nhất khi \(i-1\) bò đầu đã dẫn xong, bò \(i\) đang dẫn, đã chạy \(j\) vòng và bò \(i\) đã tốn \(k\) thể lực. Chuyển \(F[i][j+x][k+x^2]\gets F[i][j][k]+1\) và đổi bò dẫn bằng \(F[i+1][j][j]\gets F[i][j][k]\).}
\problemrow{Rebuilding Roads}{Cho một cây. Cắt một số cạnh để thu được một cây con đúng \(P\) đỉnh. Hỏi cần cắt ít nhất bao nhiêu cạnh.}{DP cây: \(F[i][j]\) là số cạnh ít nhất cần cắt trong cây con gốc \(i\) để giữ lại một cây con \(j\) đỉnh.}
\problemrow{Triangular Pastures}{Có \(N\) đoạn thẳng độ dài \(L_i\). Ghép chúng thành tam giác có diện tích lớn nhất.}{DP tập tổng: \(F[i][j]\) cho biết có thể tạo hai cạnh độ dài \(i,j\) hay không; cạnh thứ ba là tổng độ dài trừ \(i+j\). Sau đó kiểm tra điều kiện tam giác và tính diện tích.}
\problemrow{Chores}{Có \(N\) công việc; trừ công việc \(1\), mỗi công việc có một số công việc phải hoàn thành trước. Hỏi thời gian sớm nhất để hoàn thành tất cả.}{DAG DP: \(F[i]\) là thời điểm sớm nhất hoàn thành công việc \(i\), \(F[i]=t_i+\max F[j]\) trên các tiền nhiệm \(j\).}
\problemrow{Dessert}{Chèn \texttt{+}, \texttt{-} hoặc \texttt{.} giữa các số \(1,\ldots,N\) sao cho kết quả bằng \(0\).}{Liệt kê toàn bộ cách chèn toán tử.}
\problemrow{Extra Krunch}{Cho một câu, biến đổi câu sao cho không còn nguyên âm và mỗi phụ âm chỉ giữ lần xuất hiện đầu tiên.}{Mô phỏng trực tiếp.}
\problemrow{BUY LOW, BUY LOWER}{Tính số dãy con tăng dài nhất khác nhau.}{Khi có nhiều giá trị bằng nhau xuất hiện trước đó, chỉ chuyển trạng thái từ lần xuất hiện cuối cùng để tránh đếm trùng.}
\end{solutiontable}

\section{USACO 2003 Fall}

\begin{solutiontable}
\problemrow{Cow Exhibition}{Có \(N\) cặp \((A_i,B_i)\). Chọn vài cặp, đặt \(TS=\sum A_j\), \(TF=\sum B_j\). Tối đa hóa \(TS+TF\) với \(TS,TF\ge 0\).}{DP theo tổng \(TS\): \(F[i]\) là giá trị \(TF\) lớn nhất khi \(TS=i\), có dịch chỉ số để xử lý tổng âm.}
\problemrow{Milking Grid}{Cho ma trận chữ cái. Tìm ma trận con nhỏ nhất sao cho lặp vô hạn theo hai chiều sẽ chứa lại ma trận ban đầu.}{Dùng KMP tìm chu kỳ nhỏ nhất của từng hàng; bội chung nhỏ nhất của các chu kỳ là chiều rộng. Làm tương tự theo cột để có chiều cao.}
\problemrow{Popular Cows}{Trong đồ thị có hướng \(N\) đỉnh, tìm mọi đỉnh mà mọi đỉnh khác đều đi tới được.}{Đảo cạnh, co SCC. Nếu đồ thị co chỉ có một đỉnh bậc vào bằng \(0\), các đỉnh trong SCC đó là đáp án; nếu có từ hai SCC như vậy thì không có đáp án.}
\problemrow{Beauty Contest}{Tìm cặp điểm xa nhất.}{Dựng bao lồi, rồi dùng quay thước kẹp hoặc duyệt trên bao lồi.}
\problemrow{Cow Laundry}{Hàng thứ nhất \(A[i]\) nối với hàng thứ hai \(B[i]\). Mỗi lần chỉ được đổi chỗ hai phần tử kề nhau. Hỏi ít nhất bao nhiêu lần đổi để đạt \(A[i]=B[i]\).}{Quy về đếm nghịch thế sau khi ánh xạ vị trí tương ứng.}
\problemrow{Romeo Meets Juliet}{Có \(N\) con bò đang ăn cỏ trên \(P\) bãi. Tìm đoạn bãi liên tiếp dài nhất chứa không quá \(C\) con bò, một bãi có thể có nhiều bò.}{Quét hai con trỏ trên mảng số bò từng bãi.}
\problemrow{ISBN}{Chuỗi ISBN dài \(10\), chữ số cuối có thể là \(10\). Một vị trí bị mất; khôi phục vị trí đó để \(\sum i\cdot num[i]\) chia hết cho \(11\).}{Liệt kê chữ số bị mất, chú ý chữ số cuối có thể là \(10\).}
\end{solutiontable}

\section{USACO 2003 February}

\begin{solutiontable}
\problemrow{Cow Math}{Tính bội chung nhỏ nhất của ước chung lớn nhất các độ dài của mọi đường đi từ \(1\) đến \(2\).}{Tìm kiếm các đường đi. Nếu ước chung lớn nhất hiện tại đã là ước của bội chung nhỏ nhất tốt nhất thì cắt nhánh.}
\problemrow{Cow Imposters}{Có một chuỗi đích và vài chuỗi hiện có. Tìm chuỗi gần đích nhất có thể tạo từ các chuỗi hiện có bằng phép XOR.}{BFS trên các trạng thái chuỗi hoặc giá trị mã hóa của chúng.}
\problemrow{Traffic Lights}{Trên đường thẳng có đèn giao thông. Nếu gặp đèn đỏ phải dừng đến khi xanh. Mỗi giây nguyên được đổi vận tốc thêm \(-1,0,1\); khi gặp đèn đỏ vận tốc phải về \(0\). Hỏi thời gian ít nhất để đến đích với vận tốc \(0\).}{DP khả đạt \(F[i][j][k]\): ở thời điểm \(i\), vị trí \(j\), vận tốc \(k\). Chuyển sang \(k-1,k,k+1\) nếu trạng thái mới hợp lệ với đèn.}
\problemrow{Farm Tour}{Tìm hành trình ngắn nhất từ \(1\) tới \(N\) rồi từ \(N\) về \(1\), không dùng lại cạnh.}{Min-cost flow, gửi hai đơn vị luồng từ \(1\) đến \(N\).}
\problemrow{Vertical Histogram}{Thống kê số lần xuất hiện của từng chữ cái.}{Mô phỏng và in biểu đồ.}
\problemrow{Cowties}{Có \(N\) con bò, cần nối thành vòng \(1-2-\cdots-N-1\), mỗi bò đặt tại vị trí nó thích. Tối thiểu hóa độ dài vòng.}{Liệt kê vị trí của bò đầu tiên, rồi chạy DP cho các vị trí còn lại.}
\problemrow{Travel Games}{Cho một chuỗi và một từ điển. Mỗi bước có thể chèn một ký tự vào bất kỳ vị trí sao cho chuỗi mới vẫn nằm trong từ điển. Hỏi chuỗi dài nhất có thể đạt.}{Sắp từ điển theo độ dài, rồi DP từ chuỗi ngắn sang dài.}
\end{solutiontable}

\section{USACO 2003 March}

\begin{solutiontable}
\problemrow{Best Cow Fences}{Cho dãy dài \(N\). Tìm đoạn con liên tiếp có trung bình lớn nhất, độ dài ít nhất \(F\).}{Nhị phân giá trị trung bình; trừ giá trị này khỏi từng phần tử và kiểm tra có đoạn dài ít nhất \(F\) với tổng dương hay không.}
\problemrow{Cornfields}{Lưới \(N\times N\) có cao độ. Có \(K\) truy vấn hình vuông cạnh \(B\) với góc trái trên \((x,y)\), cần hiệu giữa cao nhất và thấp nhất.}{Tiền xử lý min/max cho mọi hình vuông cạnh \(B\), có thể làm thô hoặc tối ưu bằng hàng đợi đơn điệu hai chiều.}
\problemrow{Six Degrees of Cowvin Bacon}{Có \(N\) bò và \(M\) bộ phim; hai bò cùng xem một phim có khoảng cách \(1\), còn lại lấy đường đi ngắn nhất. Tìm bò có tổng khoảng cách đến các bò khác nhỏ nhất.}{Floyd--Warshall tính mọi khoảng cách, rồi duyệt tổng từng đỉnh.}
\problemrow{Herd Sums}{Đếm số cách biểu diễn \(n\) thành tổng của vài số nguyên dương liên tiếp.}{Với \(k\) số liên tiếp bắt đầu từ \(a\), \((a+a+k-1)k/2=n\). Do đó \(k\mid 2n\) và \(2n/k-k+1\) chẵn; chỉ cần duyệt \(k\).}
\problemrow{Message Decowding}{Cho một phép biến đổi chữ cái, biến đổi một câu theo bảng đó.}{Mô phỏng theo bảng ánh xạ.}
\end{solutiontable}

\section{USACO 2003 U S Open}

\begin{solutiontable}
\problemrow{Mountain Walking}{Lưới \(N\times N\) có độ cao. Đi từ \((1,1)\) đến \((N,N)\), tối thiểu hóa hiệu giữa độ cao lớn nhất và nhỏ nhất trên đường đi.}{Liệt kê độ cao thấp nhất, nhị phân độ cao cao nhất, rồi kiểm tra kết nối.}
\problemrow{Millenium Leapcow}{Lưới \(N\times N\) có số. Một quân mã chỉ được nhảy tới ô có giá trị lớn hơn. Tìm đường dài nhất; nếu nhiều nghiệm, lấy thứ tự từ điển nhỏ nhất.}{Sắp ô theo giá trị giảm dần và DP đường đi.}
\problemrow{Optimal Milking}{Có \(K\) máy và \(C\) bò, mỗi máy xử lý tối đa \(M\) bò. Tối thiểu hóa khoảng cách xa nhất mà một bò phải đi tới máy.}{Nhị phân đáp án, rồi kiểm tra bằng matching hoặc flow sau khi tính khoảng cách.}
\problemrow{Bale Figures}{Cho cách xếp các khối lập phương, tính diện tích mặt ngoài.}{Mô phỏng các mặt tiếp xúc và mặt lộ ra.}
\problemrow{Jumping Cows}{Có \(N\) số, chọn theo thứ tự để tối đa hóa tổng các số ở lượt chọn lẻ trừ tổng ở lượt chọn chẵn.}{DP \(F[i][0]\), \(F[i][1]\) là giá trị lớn nhất sau \(i\) số khi đã chọn chẵn hoặc lẻ số lượng.}
\problemrow{Lost Cows}{Khôi phục hoán vị \(1..N\) khi biết mỗi vị trí có bao nhiêu số nhỏ hơn nó đứng trước.}{Xử lý từ phải sang trái; mỗi bước lấy phần tử nhỏ thứ \(K\) trong tập còn lại.}
\problemrow{Bovine Math Geniuses}{Lặp phép biến đổi: lấy bốn chữ số giữa của số sáu chữ số, bình phương, rồi lấy sáu chữ số cuối, đến khi tạo chu trình. Tìm điểm bắt đầu, độ dài chu trình và số bước trước chu trình.}{Ghi \(F[i]\) là thời điểm lần đầu gặp số \(i\); khi gặp lại \(i\), suy ra chu trình.}
\end{solutiontable}

\section{USACO 2004 December}

\begin{solutiontable}
\problemrow{Dividing the Path}{Trên đoạn dài \(L\), đặt vòi tưới có tầm phủ \([A,B]\). Có \(N\) đoạn bắt buộc mỗi đoạn phải được một vòi phủ trọn. Hỏi ít nhất bao nhiêu vòi để phủ cả đoạn, các vòi không chồng lấn.}{DP \(F[i]\) là số vòi ít nhất để phủ đến \(i\). Chuyển \(F[i]=\min F[k]\) với \(k-i\) chẵn và \(2A\le k-i\le 2B\); duy trì min bằng hàng đợi đơn điệu.}
\problemrow{Fence Obstacle Course}{Có \(N\) hàng rào. Đi từ điểm \(S\) đến \((0,0)\), chỉ đi dọc hàng rào; tới đầu mút thì rơi thẳng xuống đến hàng rào tiếp theo. Tối thiểu hóa quãng đường ngang.}{Với mỗi hàng rào, dùng cây đoạn để biết khi rơi từ mỗi đầu mút sẽ chạm hàng rào nào; sau đó DP trên DAG theo độ cao.}
\problemrow{Cow Ski Area}{Bò trượt từ nơi cao xuống thấp, cùng độ cao cũng trượt được. Lắp thang máy giữa các ô tùy ý. Cần ít nhất bao nhiêu thang để mọi cặp ô đi tới nhau được.}{Co SCC theo khả năng trượt. Nếu chỉ còn một SCC, đáp án \(0\); ngược lại đáp án là \(\max(\#\text{nguồn},\#\text{đích})\) trong DAG co.}
\problemrow{Cleaning Shifts}{Cho \(N\) đoạn, cần ít đoạn nhất để phủ \([1,T]\).}{Sắp theo đầu trái, rồi tham lam luôn chọn đoạn có đầu phải xa nhất trong các đoạn đang dùng được.}
\problemrow{Bad Cowtractors}{Tìm cây khung lớn nhất.}{Prim hoặc Kruskal theo trọng số giảm dần.}
\problemrow{Tree Cutting}{Tìm mọi đỉnh mà nếu xóa đỉnh đó, mọi thành phần còn lại đều có số đỉnh không vượt một nửa tổng số đỉnh.}{Tính \(S[i]\) là kích thước cây con. \(F[i]=\max(\max S[\text{con}], N-S[i])\); đỉnh hợp lệ khi \(F[i]\le N/2\).}
\end{solutiontable}

\section{USACO 2004 February}

\begin{solutiontable}
\problemrow{Navigation Nightmare}{Cây trên mặt phẳng. Truy vấn \((F_1,F_2,I)\): dùng chỉ \(I\) cạnh đầu, hỏi khoảng cách Manhattan giữa \(F_1,F_2\), hoặc \(-1\) nếu chưa liên thông.}{Xây cả cây để tính tọa độ và khoảng cách, rồi với mỗi truy vấn kiểm tra liên thông khi chỉ dùng \(I\) cạnh đầu bằng DSU offline hoặc xử lý theo thứ tự \(I\).}
\problemrow{Cow Marathon}{Tìm đường dài nhất trên cây.}{Hai lần BFS/DFS: từ đỉnh bất kỳ tìm đầu mút xa nhất, rồi từ đó tìm khoảng cách xa nhất.}
\problemrow{Distance Queries}{Truy vấn khoảng cách giữa hai đỉnh trên cây.}{Gọi \(D[i]\) là khoảng cách từ \(i\) tới gốc, \(k=\operatorname{lca}(i,j)\). Khi đó \(\operatorname{dist}(i,j)=D[i]+D[j]-2D[k]\).}
\problemrow{Distance Statistics}{Đếm số cặp đỉnh trên cây có khoảng cách nhỏ hơn \(M\).}{Dùng cây cân bằng/DSU on tree theo khóa \(d[i]\), khoảng cách từ gốc. Khi xử lý đỉnh \(u\), giữ cấu trúc của con nặng, truy vấn các con nhẹ rồi trộn vào để đếm cặp hợp lệ.}
\end{solutiontable}

\section{USACO 2004 March}

\begin{solutiontable}
\problemrow{Moo University -- Team Tryouts}{Có \(N\) cặp \((H_i,W_i)\). Chọn nhiều cặp nhất sao cho với \(h,w\) là min của các cặp đã chọn, mọi cặp thỏa \(A(H_i-h)+B(W_i-w)\le C\).}{Biến đổi thành \(AH_i+BW_i-C\le Ah+Bw\). Sắp theo khóa này, liệt kê \(h,w\), dùng heap để duy trì và thống kê các ứng viên.}
\problemrow{Moo University -- Emergency Pizza Order}{Có \(C\) người, \(T\) loại đồ. Mỗi người thích một số loại. Cấp cho mỗi người \(K\) loại họ thích, và không có hai người nhận y hệt nhau. Tối đa hóa số người được thỏa.}{Quy về matching.}
\problemrow{Moo University -- Financial Aid}{Có \(C\) cặp \((A_i,B_i)\). Chọn \(N\) cặp, \(N\) lẻ, sao cho trung vị của \(A\) lớn nhất và tổng \(B\le F\).}{Liệt kê trung vị. Với phía nhỏ hơn và lớn hơn trung vị, lấy \(N/2\) giá trị \(B\) nhỏ nhất bằng heap; kiểm tra tổng ngân sách.}
\end{solutiontable}

\section{USACO 2004 November}

\begin{solutiontable}
\problemrow{Apple Catching}{Có hai vị trí; trong \(T\) giây, mỗi giây một quả táo rơi ở một vị trí. Bạn bắt táo và được di chuyển tối đa \(W\) lần. Hỏi bắt được nhiều nhất bao nhiêu quả.}{DP \(F[i][j]\): tại thời điểm \(i\), đã di chuyển \(j\) lần, số táo tối đa. Vị trí hiện tại là \((j+1)\bmod 2\).}
\problemrow{Lake Counting}{Cho bản đồ \(N\times M\), đếm số thành phần liên thông của nước.}{BFS/DFS hoặc DSU.}
\problemrow{Til the Cows Come Home}{Tìm đường đi ngắn nhất từ \(1\) đến \(N\).}{Dijkstra hoặc SPFA.}
\problemrow{Who's in the Middle}{Tìm trung vị của một dãy.}{Sắp xếp rồi in phần tử giữa.}
\problemrow{Bull Math}{Nhân hai số lớn.}{Nhân số nguyên lớn theo từng chữ số.}
\problemrow{Bank Interest}{Có \(M\) tiền, mỗi năm nhận lãi \(R\%\). Hỏi sau \(Y\) năm có bao nhiêu tiền.}{Mô phỏng từng năm với làm tròn theo yêu cầu đề.}
\end{solutiontable}

\section{USACO 2004 U S Open}

\begin{solutiontable}
\problemrow{Cube Stacking}{Ban đầu có \(N\) khối lập phương. Hai thao tác: đặt cả chồng chứa \(X\) lên chồng chứa \(Y\), hoặc hỏi dưới \(X\) có bao nhiêu khối.}{DSU có trọng số lưu cha của chồng và số khối nằm dưới mỗi phần tử.}
\problemrow{The Cow Lineup}{Cho một dãy. Tìm độ dài ngắn nhất của một dãy không phải dãy con của dãy gốc.}{Quét từ đầu đến khi gặp đủ \(1..K\), cắt một đoạn rồi quét tiếp. Nếu cắt được \(M\) đoạn thì đáp án là \(M+1\).}
\problemrow{MooFest}{Có \(N\) bò ở vị trí \(X_i\), âm lượng \(V_i\). Năng lượng giữa \(i,j\) là \(|X_i-X_j|\max(V_i,V_j)\). Tính tổng trên mọi cặp.}{Sắp theo \(V_i\), thêm dần vị trí vào cây Fenwick/cây đoạn để tính tổng khoảng cách tới các bò đã thêm.}
\problemrow{Turning in Homework}{Giáo viên thu bài ở các lớp vị trí \(P_i\), bài lớp \(i\) chỉ thu sau thời điểm \(T_i\). Tốc độ \(1\); cần thu hết và đến vị trí \(B\) sớm nhất.}{Sắp theo \(P_i\). DP đoạn \(F[i][j][0/1]\): các lớp ngoài đoạn \([i,j]\) đã thu, hiện ở \(i\) hoặc \(j\), thời gian nhỏ nhất.}
\end{solutiontable}

\section{USACO 2005 December}

\begin{solutiontable}
\problemrow{Alignment of the Planets}{Trên mặt phẳng có \(N\) điểm, đếm số bộ ba điểm thẳng hàng.}{Liệt kê các bộ ba hoặc chuẩn hóa hệ số góc theo từng điểm tùy giới hạn.}
\problemrow{Finding Boving Roots}{Tìm số nhỏ nhất mà căn bậc hai của nó có \(L\) chữ số sau dấu thập phân trùng với số đã cho.}{Liệt kê phần nguyên, bình phương sau khi làm tròn thành số nguyên, rồi kiểm tra phần thập phân.}
\problemrow{Cow Bowling}{Tìm đường từ đỉnh xuống đáy trong tam giác số sao cho tổng lớn nhất.}{DP tam giác kinh điển từ trên xuống hoặc dưới lên.}
\problemrow{Cow Patterns}{Hai dãy được xem như nhau nếu cùng độ dài và quan hệ \(<,=,>\) giữa mọi cặp vị trí tương ứng giống nhau. Đếm số đoạn con của một dãy giống dãy mẫu.}{KMP trên quan hệ thứ hạng. Hai phần tử khớp khi số phần tử nhỏ hơn, bằng và lớn hơn trong tiền tố tương ứng giống nhau.}
\problemrow{Barn Expansion}{Có \(N\) hình chữ nhật không giao nhau nhưng có thể chạm nhau. Đếm số hình có thể mở rộng đồng thời cả bốn cạnh một đoạn dương.}{Vì không giao nhau, sắp theo \(x\) để kiểm tra cản trở trên các đường thẳng đứng, rồi sắp theo \(y\) làm tương tự.}
\problemrow{Layout}{Trên đường thẳng có \(N\) điểm. Có ràng buộc khoảng cách \(A,B\) không quá \(D\), ràng buộc không nhỏ hơn \(D\), và thứ tự điểm không đảo. Tìm khoảng cách lớn nhất có thể giữa \(1\) và \(N\).}{Mô hình chuẩn là ràng buộc hiệu. Cách khác: duy trì khoảng khả thi \([l_i,r_i]\), nới chặt bằng từng loại quan hệ đến khi ổn định hoặc phát hiện vô nghiệm.}
\problemrow{Knights of Ni}{Bản đồ \(N\times M\), \(2\) là xuất phát, \(3\) là đích. Cần đi từ \(2\) tới một ô \(4\), rồi tới \(3\), với ít bước nhất.}{BFS một lần trên trạng thái gồm ô hiện tại và cờ đã ghé ô \(4\) hay chưa.}
\problemrow{Cleaning Shifts}{Có \(N\) bò; thuê bò \(j\) làm từ \(T1_j\) đến \(T2_j\) với giá \(S_j\). Phủ mọi thời điểm từ \(M\) đến \(E\) với chi phí nhỏ nhất.}{DP \(F[i]\) là chi phí nhỏ nhất để phủ đến thời điểm \(i\). Với đoạn kết thúc ở \(i\), \(F[i]=\min(F[t]+S_j)\) trên \(t\in[T1_j,T2_j-1]\); dùng cây đoạn để lấy min.}
\problemrow{Scales}{Có các quả cân, tìm trọng lượng lớn nhất không vượt \(C\) có thể tạo được. Giả sử \(W_i+W_{i+1}\le W_{i+2}\).}{Tìm kiếm có cắt nhánh. Nếu \(M+W_k+W_{k-1}\le C\) thì bắt buộc lấy \(W_k\); nếu \(M+W_k\le C<M+W_k+W_{k-1}\) thì chỉ một trong hai quả \(k,k-1\) có thể được lấy.}
\end{solutiontable}

\section{USACO 2005 February}

\begin{solutiontable}
\problemrow{Jersey Politics}{Có \(3K\) thành phố, chia thành \(3\) nhóm mỗi nhóm \(K\) thành phố sao cho có \(2\) nhóm có tổng dân số lớn hơn \(500K\).}{Lấy \(2K\) thành phố lớn nhất, chia tùy ý thành hai nhóm rồi điều chỉnh ngẫu nhiên hoặc hoán đổi cục bộ.}
\problemrow{Secret Milking Machine}{Có \(N\) thành phố, \(M\) cạnh. Cần \(T\) đường đi từ \(1\) đến \(N\), không dùng chung cạnh, sao cho cạnh dài nhất được dùng là nhỏ nhất.}{Nhị phân độ dài cạnh tối đa, giữ các cạnh không vượt ngưỡng rồi kiểm tra có \(T\) đường cạnh-rời bằng max flow.}
\problemrow{Aggressive cows}{Trên trục \(x\) có \(N\) điểm, chọn \(C\) điểm để khoảng cách nhỏ nhất giữa hai điểm được chọn là lớn nhất.}{Nhị phân khoảng cách, kiểm tra tham lam bằng cách đặt bò càng xa càng tốt từ trái sang phải.}
\problemrow{Part Acquisition}{Có \(K\) đỉnh và \(N\) cạnh có hướng, tìm đường ngắn nhất từ \(1\) đến \(K\).}{BFS nếu cạnh không trọng số.}
\problemrow{Rigging the Bovine Election}{Chọn \(7\) ô liên thông trong lưới \(5\times5\), số ô loại \(J\) nhiều hơn loại \(H\). Hỏi có bao nhiêu cách.}{Tìm kiếm chọn/không chọn từng ô; khi đủ \(7\) ô thì kiểm tra liên thông và số lượng \(J,H\).}
\end{solutiontable}

\section{USACO 2005 January}

\begin{solutiontable}
\problemrow{Muddy Fields}{Trong ma trận \(R\times C\), vài ô là \texttt{*}. Dùng ít tấm ván nhất để phủ mọi \texttt{*}, không phủ ô \texttt{.}.}{Đánh số các đoạn \texttt{*} liên tiếp theo hàng và theo cột, rồi làm matching hai phía.}
\problemrow{The Wedding Juicer}{Lưới \(N\times M\) có độ cao. Tính lượng nước đọng sau mưa lớn.}{Đưa toàn bộ biên vào heap min; mỗi lần lấy ô thấp nhất, lan sang ô chưa thăm. Nếu ô kề thấp hơn mực hiện tại thì cộng nước bằng hiệu độ cao và đẩy với mực mới.}
\problemrow{Naptime}{Có \(N\) số trên vòng tròn, chọn \(B\) số. Điểm đầu tiên của mỗi đoạn được chọn liên tiếp không được tính. Tối đa hóa điểm.}{DP vòng. \(F[i][j][0/1]\) là điểm tối đa xét \(i\) số đầu, đã chọn \(j\) số, trạng thái số \(i\). Tách hai trường hợp số đầu có được tính điểm hay không.}
\problemrow{Sumsets}{Đếm số cách viết \(N\) thành tổng của các lũy thừa của \(2\).}{\(F[i]\) là số cách của \(i\). Nếu \(i\) lẻ, \(F[i]=F[i-1]\); nếu \(i\) chẵn, \(F[i]=F[i-1]+F[i/2]\).}
\problemrow{Watchcow}{Đồ thị vô hướng, mỗi cạnh được đi hai lần, mỗi chiều một lần. In một hành trình Euler phù hợp.}{DFS theo cạnh có hướng tương ứng với hai chiều của mỗi cạnh vô hướng.}
\problemrow{Moo Volume}{Trên trục số có \(N\) điểm, tính tổng khoảng cách có hướng giữa mọi cặp rồi nhân đôi theo yêu cầu.}{Sắp \(d_i\). Đóng góp của điểm \(i\) với các điểm trước là \((i-1)d_i-\sum_{j<i}d_j\); cuối cùng nhân \(2\).}
\end{solutiontable}

\section{USACO 2005 March}

\begin{solutiontable}
\problemrow{Ombrophobic Bovines}{Có \(F\) chuồng, ban đầu chuồng \(i\) có \(A_i\) bò và chứa tối đa \(B_i\). Chuyển bò giữa chuồng để tối thiểu hóa khoảng cách xa nhất một con bò phải đi.}{Tính mọi đường ngắn nhất, nhị phân đáp án. Với ngưỡng \(D\), tách mỗi chuồng thành \(i,i'\), nguồn nối \(i\) dung lượng \(A_i\), \(i'\) nối đích dung lượng \(B_i\), và \(i\to j'\) vô hạn nếu khoảng cách \(i,j\le D\).}
\problemrow{Space Elevator}{Có \(K\) loại gạch, loại \(i\) cao \(H_i\), có \(C_i\) viên, và không được vượt quá cao độ \(A_i\). Tìm độ cao lớn nhất.}{Sắp theo \(A_i\). DP khả đạt \(Can[i][j]\) bằng \(i\) loại đầu; chuyển bằng cách thử số viên loại tiếp theo.}
\problemrow{Yogurt factory}{Trong \(N\) ngày, chi phí sản xuất sữa ngày \(i\) là \(C_i\), nhu cầu \(Y_i\). Sữa dư có thể lưu kho với phí \(S\) mỗi đơn vị mỗi ngày. Tối thiểu hóa chi phí.}{Tham lam giữ ngày sản xuất hiệu quả nhất đến hiện tại: chi phí dùng cho ngày \(i\) là min giữa sản xuất hôm nay và sản xuất từ ngày tốt nhất trước đó cộng phí lưu kho.}
\problemrow{Checking an Alibi}{Có \(F\) chuồng, \(P\) cạnh, và \(C\) con bò ở các chuồng khác nhau. Tìm những bò có thể tới chuồng \(1\) trong thời gian \(M\).}{Chạy đường ngắn nhất từ \(1\), rồi kiểm tra khoảng cách của từng bò.}
\problemrow{Out of hay}{Có \(N\) đỉnh, \(M\) cạnh. Giữ một số cạnh sao cho đồ thị liên thông và cạnh lớn nhất nhỏ nhất.}{Bản chất là trọng số lớn nhất trong MST; có thể nhị phân đáp án rồi BFS kiểm tra liên thông.}
\end{solutiontable}

\section{USACO 2005 November}

\begin{solutiontable}
\problemrow{Securing the Barn}{Có \(N\) chữ cái, chọn \(M\) chữ cái, sắp tăng thành chuỗi, trong đó có ít nhất một nguyên âm.}{Duyệt mọi tổ hợp.}
\problemrow{Hopscotch}{Ma trận \(5\times5\). Bắt đầu từ ô bất kỳ, đi sang ô kề \(5\) bước để tạo số \(6\) chữ số, được phép lặp ô. Đếm số khác nhau.}{DFS từ từng ô, dùng hash/set để lưu số tạo được.}
\problemrow{Satellite Photographs}{Trong bản đồ \(N\times M\), tìm thành phần liên thông chứa nhiều ký tự \texttt{*} nhất.}{BFS/DFS từng thành phần.}
\problemrow{Asteroids}{Lưới \(N\times N\) có vật thể. Mỗi lần xóa toàn bộ vật thể trên một hàng hoặc cột. Cần ít nhất bao nhiêu lần để xóa hết.}{Bài toán phủ đỉnh nhỏ nhất trong đồ thị hai phía, giải bằng matching cực đại.}
\problemrow{Grazing on the Run}{Trên trục số có \(N\) điểm, xuất phát tại \(L\), tốc độ \(1\). Nếu đến điểm \(i\) tại thời điểm \(T_i\), tối thiểu hóa \(\sum T_i\).}{DP đoạn \(F[i][j][0/1]\): các điểm trong đoạn \([i,j]\) đã đi qua, hiện ở \(i\) hoặc \(j\), chi phí nhỏ nhất.}
\problemrow{Walk the Talk}{Ma trận ký tự \(N\times M\) và một từ điển. Từ vị trí bất kỳ, mỗi bước chỉ nhảy lên phải. Đếm số chuỗi trong từ điển có thể tạo được.}{Xử lý từng từ. DP \(F[l][i][j]\) là số cách sau \(l\) bước đang ở ô \((i,j)\); chuyển đến các ô phía trên bên phải.}
\problemrow{City Skyline}{Có dải chiều rộng \(W\), tại các vị trí \(W_i\) có độ cao \(H_i\). Đếm số hình chữ nhật cực đại không thể mở rộng.}{Với mỗi \(h_i\), tính biên trái/phải xa nhất mà \(h_i\) vẫn là chiều cao thấp nhất; sắp và loại trùng các cặp biên.}
\problemrow{Cow Acrobats}{Mỗi bò có trọng lượng \(W_i\) và sức chịu \(S_i\). Rủi ro của một bò là tổng trọng lượng các bò phía trên trừ \(S_i\). Tối thiểu hóa rủi ro lớn nhất.}{Tham lam sắp theo \(W_i+S_i\).}
\problemrow{Ant Counting}{Có \(N\) số thuộc \(1..T\). Đếm số tập con có kích thước từ \(A\) đến \(B\).}{DP \(F[i][j]\): dùng các số \(1..i\), tạo kích thước \(j\). Chuyển \(F[i][j]=\sum_k F[i-1][j-k]\) với \(k\) không vượt số lần xuất hiện của \(i\).}
\end{solutiontable}

\section{USACO 2005 October}

\begin{solutiontable}
\problemrow{Bovine Birthday}{Biết ngày 1/1/1990 là thứ Hai, hỏi một ngày cho trước là thứ mấy.}{Mô phỏng lịch, xét năm nhuận.}
\problemrow{Max Factor}{Tìm số có thừa số nguyên tố lớn nhất trong một tập số.}{Phân tích thừa số nguyên tố bằng chia thử.}
\problemrow{Skiing}{Lưới \(N\times M\) có độ cao. Đi từ \((1,1)\) tới \((N,M)\); vận tốc đổi thành \(v\cdot2^{h_1-h_2}\) khi đi giữa hai ô. Tìm thời gian ít nhất.}{Sau khi tới một ô, vận tốc được xác định bởi độ cao, nên dùng SPFA/đường ngắn nhất trên trạng thái ô.}
\problemrow{Flying Right}{Trên đường thẳng có \(N\) nông trại và \(K\) nhóm bò cần đi giữa hai nông trại. Máy bay chứa tối đa \(C\) bò, bay từ \(1\) đến \(N\) rồi từ \(N\) về \(1\). Tối đa hóa số bò được chở.}{Làm riêng hai chiều. Sắp theo đầu phải, tham lam nhận nhóm nếu còn sức chứa trên toàn đoạn; dùng cây đoạn để mô phỏng tải.}
\problemrow{Close Encounter}{Cho một phân số, tìm phân số tối giản khác nó và gần nó nhất, tử và mẫu đều nhỏ hơn \(32768\).}{Liệt kê mẫu, tính hai tử gần nhất và kiểm tra tối giản.}
\problemrow{Allowance}{Có \(N\) mệnh giá \(V_i\), mỗi loại \(B_i\) tờ; sau khi sắp, \(V_{i+1}\) là bội của \(V_i\). Tối đa hóa số lần trả tiền, mỗi lần tổng mệnh giá ít nhất \(C\).}{Tham lam từ mệnh giá lớn xuống dùng nhiều nhất nhưng không vượt \(C\); phần thiếu dùng tờ nhỏ nhất vượt ngưỡng còn lại để bù.}
\end{solutiontable}

\section{USACO 2005 U S Open}

\begin{solutiontable}
\problemrow{Lazy Cows}{Trong hình chữ nhật \(2\times B\) có \(N\) ô có bò. Dùng \(K\) hình chữ nhật phủ chúng, tối thiểu hóa tổng diện tích.}{DP nén trạng thái: \(F[i][j][0..3]\) là diện tích nhỏ nhất sau \(i\) bò, đã dùng \(j\) hình; trạng thái mô tả phủ hàng trên, hàng dưới, cả hai hàng chung, hoặc tách riêng.}
\problemrow{Expedition}{Trên đường thẳng có \(N\) trạm xăng; trạm cách thành phố \(D_i\), có \(E_i\) xăng. Bạn cách thành phố \(L\), có \(P\) xăng. Hỏi ít nhất dừng bao nhiêu trạm.}{Tham lam với heap max: khi không tới được trạm tiếp theo, lấy lượng xăng lớn nhất trong các trạm đã đi qua.}
\problemrow{Around the world}{Trên đường tròn có \(N\) điểm và \(M\) cạnh; cạnh giữa hai điểm là cung ngắn hơn \(180^\circ\). Tìm ít cạnh nhất để đi từ điểm \(1\) một vòng rồi về lại chính nó.}{BFS trên \(F[i][j]\): đang ở điểm \(i\), đã quấn quanh \(j\) vòng; \(j<0\) biểu thị chiều ngược kim đồng hồ. Khi đi qua điểm \(1\), cập nhật số vòng.}
\problemrow{Landscaping}{Có một dãy núi. Cần phá ít đá nhất để số đỉnh núi không quá \(K\).}{Tham lam mỗi lần phá đỉnh cần ít đá nhất, cập nhật lân cận.}
\problemrow{Waves}{Một số viên đá được ném xuống nước tại các thời điểm và vị trí khác nhau. Hỏi hình dạng sóng tại thời điểm \(R\).}{Mô phỏng đóng góp của từng viên đá.}
\problemrow{Navigating the City}{Bản đồ \((2N-1)\times(2M-1)\). Đi từ \(S\) đến \(E\), chỉ qua ô \texttt{+} và phải có \texttt{-} hoặc \texttt{|} nối giữa hai giao điểm. In đường ngắn nhất.}{BFS trên các giao điểm hợp lệ, lưu cha để khôi phục đường đi.}
\problemrow{Disease Management}{Có \(D\) loại bệnh. Chọn nhiều bò nhất sao cho tổng số loại bệnh xuất hiện trong nhóm không vượt \(K\).}{Liệt kê tập bệnh được phép, đếm số bò có tập bệnh là con của tập đó.}
\problemrow{Muddy roads}{Có \(N\) đoạn, dùng ván dài \(L\) phủ toàn bộ các đoạn. Hỏi cần ít ván nhất.}{Sắp đoạn theo đầu trái, đặt ván tham lam từ điểm chưa phủ đầu tiên.}
\end{solutiontable}

\section{USACO 2006 December}

\begin{solutiontable}
\problemrow{Parkside's Triangle}{Điền số bắt đầu từ \(S\), theo từng cột trong tam giác.}{Mô phỏng quy luật điền.}
\problemrow{Wormholes}{Có \(N\) nông trại, \(M\) đường dương hai chiều và \(W\) đường âm một chiều. Hỏi có chu trình âm hay không.}{Bellman--Ford hoặc SPFA phát hiện chu trình âm.}
\problemrow{The Fewest Coins}{Có \(N\) loại tiền, giá trị \(V_i\), có \(C_i\) tờ. Mua món giá \(T\), tối thiểu hóa số tờ đưa ra cộng số tờ nhận lại.}{DP có giới hạn cho số tờ tối thiểu để trả \(i\), DP không giới hạn cho số tờ tối thiểu ông chủ trả lại \(j\); liệt kê tổng tiền đã đưa.}
\problemrow{Milk Patterns}{Dãy dài \(N\). Tìm đoạn con dài nhất xuất hiện ít nhất \(K\) lần, có thể chồng lấn.}{Nhị phân độ dài, kiểm tra bằng hash hoặc suffix array.}
\problemrow{Cow Picnic}{Có \(N\) nông trại, \(M\) cạnh có hướng, \(K\) bò đứng ở vài nông trại. Đếm số nông trại mà cả \(K\) bò đều đi tới được.}{Từ vị trí mỗi bò chạy BFS/DFS, cộng bộ đếm số lần tới từng đỉnh.}
\problemrow{Cow Roller Coaster}{Có \(N\) đoạn \([X_i,X_i+W_i)\), chi phí \(C_i\), độ vui \(F_i\). Phủ đúng \([0,L)\) bằng các đoạn không chồng lấn, tổng chi phí không quá \(B\), tối đa hóa độ vui.}{DP \(G[i][j]\): đoạn \([0,i)\) đã phủ, chi phí \(j\), độ vui lớn nhất. Với đoạn kết thúc tại \(i\), chuyển từ \(G[X_k][j-C_k]+F_k\).}
\problemrow{River Hopscotch}{Đích cách xuất phát \(L\), giữa đường có \(N\) tảng đá. Loại \(M\) tảng để khoảng nhảy ngắn nhất trên đường là lớn nhất.}{Nhị phân khoảng cách nhỏ nhất, kiểm tra tham lam số đá cần loại.}
\end{solutiontable}

\section{USACO 2006 February}

\begin{solutiontable}
\problemrow{The Moronic Cowmpouter}{Chuyển một số sang hệ cơ số \(-2\).}{Tương tự đổi sang cơ số \(2\): nếu số hiện tại lẻ thì chữ số là \(1\), sau đó trừ chữ số và chia cho \(-2\).}
\problemrow{DNA Assembly}{Có \(N\) chuỗi, tìm chuỗi ngắn nhất chứa tất cả chúng làm chuỗi con liên tiếp.}{Liệt kê thứ tự ghép các chuỗi, dùng phần chồng lấn lớn nhất giữa hai chuỗi liên tiếp, lấy phương án ngắn nhất.}
\problemrow{Cow Phrasebook}{Có \(M\) chuỗi mẫu; trong \(N\) chuỗi tiếp theo, đếm bao nhiêu chuỗi là tiền tố của một chuỗi mẫu.}{Xây trie từ \(M\) chuỗi mẫu; với mỗi chuỗi truy vấn, đi trên trie để kiểm tra tiền tố.}
\problemrow{Cellphones}{Có \(B\) nút và \(L\) chữ cái đầu của bảng chữ cái, cùng \(D\) từ. Phân chữ cái vào nút để tối đa hóa số từ có cùng chuỗi nút.}{Tìm kiếm cách gán chữ cái cho từng nút, sau đó mã hóa mọi từ thành chuỗi nút và đếm tần suất lớn nhất.}
\problemrow{Steady Cow Assignment}{Có \(N\) bò và \(B\) chuồng, mỗi chuồng có sức chứa. Mỗi bò xếp hạng các chuồng. Gán bò vào chuồng để tối thiểu hóa độ rộng hạng được dùng.}{Liệt kê hoặc nhị phân khoảng hạng \([l,r]\); mỗi lần kiểm tra bằng matching/flow có sức chứa.}
\problemrow{Treats for the Cows}{Dãy số có trọng số. Mỗi lần lấy một số ở đầu hoặc cuối; số lấy thứ \(i\) đóng góp trọng số nhân \(i\). Tối đa hóa tổng.}{DP \(F[i][j]\): đã lấy \(i\) số bên trái và \(j\) số bên phải, giá trị lớn nhất.}
\problemrow{Backward Digit Sums}{Tìm hoán vị \(1..N\) sao cho tổng có trọng số nhị thức bằng \(S\), in hoán vị thứ tự từ điển nhỏ nhất.}{Tìm kiếm theo thứ tự từ điển, dùng trọng số Pascal để tính tổng.}
\problemrow{Stall Reservations}{Cho \(N\) đoạn, chia thành ít nhóm nhất sao cho trong mỗi nhóm các đoạn không chồng lấn.}{Sắp theo đầu trái; dùng heap theo đầu phải nhỏ nhất của mỗi nhóm, tái sử dụng nhóm nếu không giao, nếu không mở nhóm mới.}
\end{solutiontable}

\section{USACO 2006 January}

\begin{solutiontable}
\problemrow{Stump Removal}{Có \(N\) gốc cây cao \(H_i\). Khi kích nổ một gốc, các gốc thấp hơn ở hai bên cũng đổ lan. Hỏi ít nhất cần kích nổ những gốc nào và in chỉ số.}{Quét từ trái sang phải; nếu bên phải hiện tại có gốc thấp hơn hoặc bằng, kích nổ vị trí này rồi lan hết phần có thể đổ.}
\problemrow{Finicky Grazers}{Có \(N\) bò trên đoạn \([0,L]\). Di chuyển bò để khoảng cách hai bò kề là \(D\) hoặc \(D+1\), với tổng quãng di chuyển nhỏ nhất.}{Đặt \(D=N/(L-1)\). DP \(O(NL)\) có thể giảm vì vị trí bò \(i+1\) chỉ nằm trong khoảng \(iD\) đến \(i(D+1)\).}
\problemrow{The Water Bowls}{Có \(20\) bát, \(0\) là ngửa, \(1\) là úp. Một thao tác đổi trạng thái ba bát kề nhau. Hỏi ít nhất bao nhiêu thao tác để tất cả thành \(0\).}{BFS trên trạng thái bitmask \(20\) bit.}
\problemrow{Redundant Paths}{Cần thêm ít nhất bao nhiêu cạnh để giữa mọi cặp đỉnh có hai đường đi cạnh-rời.}{Co các thành phần không chứa cầu thành một khối, thu được cây cầu. Đáp án \((\#\text{lá}+1)/2\).}
\problemrow{Roping the Field}{Có \(N\) điểm tạo thành đa giác lồi và \(G\) hình tròn bán kính \(R\). Chọn nhiều dây chéo nhất sao cho dây không cắt hình tròn nào và các dây không cắt nhau.}{DP khoảng \(F[i][j]\): số dây tối đa giữa hai điểm \(i,j\). \(F[i][j]=\max_k(F[i][k]+F[k][j])+V[i][j]\), với \(V[i][j]=1\) nếu dây \(ij\) hợp lệ.}
\problemrow{Corral the Cows}{Trên mặt phẳng có \(N\) điểm. Tìm hình vuông nhỏ nhất chứa ít nhất \(C\) điểm.}{Trong tối ưu có hai điểm nằm trên hai cạnh. Liệt kê hai điểm làm đường song song với cạnh hình vuông, rồi tìm cạnh nhỏ nhất chứa đủ \(C\) điểm.}
\problemrow{The Cow Prom}{Đếm số nhóm bò đi prom theo quan hệ đi được qua lại.}{Co SCC và đếm các SCC có hơn một đỉnh.}
\problemrow{Dollar Dayz}{Viết \(N\) thành tổng các số nguyên dương không vượt \(K\), đếm số cách.}{Truy hồi đơn giản: \(F[i][j]=F[i][j-1]+F[i-j][j]\).}
\problemrow{The Grove}{Từ \texttt{*}, đi một vòng trở lại \texttt{*} sao cho bao quanh toàn bộ các ô \texttt{X}. Tìm số bước ít nhất.}{Lấy ô \texttt{X} trên cùng, trái nhất làm mốc. Lần đi thứ nhất từ xuất phát sang phải tới mốc, lần thứ hai từ mốc sang trái về xuất phát; dùng BFS theo ràng buộc hướng.}
\end{solutiontable}

\section{USACO 2006 November}

\begin{solutiontable}
\problemrow{Fence Repair}{Có \(N\) thanh gỗ dài \(L_i\), cần cắt từ một thanh tổng độ dài \(\sum L_i\). Mỗi lần cắt tốn chi phí bằng độ dài thanh bị cắt. Tối thiểu hóa chi phí.}{Làm ngược: mỗi lần gộp hai thanh, chi phí là tổng độ dài của chúng. Tham lam Huffman, luôn gộp hai thanh ngắn nhất.}
\problemrow{Corn Fields}{Ruộng \(N\times M\), ô \(1\) có thể trồng. Các ô được chọn không được kề cạnh. Đếm số cách chọn.}{DP nén trạng thái theo hàng: \(F[i][s]\) là số cách đến hàng \(i\) với trạng thái hàng \(i\) là \(s\), kiểm tra hợp lệ trong hàng và giữa hai hàng.}
\problemrow{Roadblocks}{Tìm đường đi ngắn thứ hai từ \(1\) đến \(N\).}{Có thể tính đường ngắn nhất rồi thử bỏ từng cạnh trên đường đó và chạy lại; cách chuẩn hơn là Dijkstra lưu hai khoảng cách tốt nhất mỗi đỉnh.}
\problemrow{Bad Hair Day}{\(N\) bò đứng thành hàng, bò \(i\) cao \(H_i\). Nếu \(i<j\), \(H_i>H_j\), và giữa chúng không có bò cao hơn \(H_i\), thì \(i\) thấy \(j\). Tính tổng số bò nhìn thấy.}{Tìm với mỗi bò vị trí đầu tiên bên phải có chiều cao ít nhất bằng nó; có thể dùng stack đơn điệu hoặc cây đoạn.}
\problemrow{Big Square}{Trên ruộng \(N\times N\) có một số ô \texttt{J} và \texttt{B}. Có thể đặt thêm một \texttt{J}. Tìm hình vuông lớn nhất có bốn đỉnh đều là \texttt{J}, cạnh không nhất thiết song song trục.}{Liệt kê hai đỉnh của một đường chéo, suy ra hai đỉnh còn lại và kiểm tra hợp lệ, cho phép một đỉnh cần thêm \texttt{J}.}
\problemrow{Round Numbers}{Đếm số Round Numbers trong \([A,B]\), tức khi viết nhị phân có số bit \(0\) không ít hơn số bit \(1\).}{Tính hàm đếm trên \([1,N]\) bằng digit DP theo nhị phân, rồi lấy \(count(B)-count(A-1)\).}
\end{solutiontable}

\section{USACO 2007 December}

\begin{solutiontable}
\problemrow{Bookshelf}{Có \(N\) bò cao \(H_i\). Cần ít bò nhất sao cho tổng chiều cao ít nhất \(B\).}{Sắp chiều cao giảm dần, lấy từ cao xuống thấp.}
\problemrow{Bookshelf2}{Có \(N\) bò cao \(H_i\). Tìm tổng chiều cao nhỏ nhất vượt quá \(B\).}{Tìm kiếm tập con hoặc DP subset-sum tùy giới hạn.}
\problemrow{Card Stacking}{Có \(N\) bò và \(K\) lá bài, trong đó \(M=K/N\) lá tốt. \(N\) bò ngồi vòng tròn. Hỏi đặt các lá tốt ở vị trí nào để chúng đều được chia cho người chia bài.}{Mô phỏng quá trình chia bài theo vòng.}
\problemrow{Sightseeing Cows}{Có \(L\) đỉnh, \(P\) cạnh có hướng; đỉnh có giá trị \(F_i\), cạnh có chi phí \(T_i\). Tìm chu trình tối đa hóa \(\sum F_i/\sum T_i\).}{Nhị phân tỉ lệ \(x\); đặt trọng số cạnh thành \(T_i x-F_i\) và kiểm tra có chu trình âm.}
\problemrow{Gourmet Grazers}{Bò \(i\) cần thức ăn giá ít nhất \(A_i\) và độ mềm ít nhất \(B_i\). Có \(M\) phần thức ăn giá \(C_i\), độ mềm \(D_i\). Tối thiểu hóa tổng giá để thỏa mọi bò.}{Sắp thức ăn theo \(C_i\) tăng dần; với mỗi phần, gán cho bò chưa thỏa có \(A_i\le C_i\) và yêu cầu \(B_i\) lớn nhất nhưng không vượt \(D_i\).}
\problemrow{Best Cow Line, Gold}{Chuỗi dài \(N\). Mỗi lần lấy ký tự ở đầu hoặc cuối để tạo chuỗi kết quả nhỏ nhất theo thứ tự từ điển.}{Tham lam: nếu hai đầu khác nhau, lấy ký tự nhỏ hơn; nếu bằng nhau, so sánh sâu vào trong. Bản Gold dùng suffix array để tiền xử lý so sánh.}
\problemrow{Charm Bracelet}{Có \(N\) viên đá, mỗi viên nặng \(W_i\) và tăng sức hấp dẫn \(D_i\). Tối đa hóa tổng \(D_i\) với trọng lượng không quá \(M\).}{Knapsack \(0/1\).}
\problemrow{Building Roads}{Có \(N\) nông trại, một số cặp đã có đường. Cần thêm đường ngắn nhất để toàn bộ liên thông.}{Cạnh đã có đặt trọng số \(0\), các cạnh còn lại theo khoảng cách Euclid; chạy MST.}
\problemrow{Mud Puddles}{Đi từ \((0,0)\) tới vị trí ký hiệu \texttt{B}; có \(N\) ô không đi được.}{BFS trên lưới.}
\end{solutiontable}

\section{USACO 2007 February}

\begin{solutiontable}
\problemrow{Building A New Barn}{Có \(N\) chuồng. Thêm một chuồng mới để tổng khoảng cách Manhattan tới các chuồng cũ nhỏ nhất, nhưng không trùng vị trí chuồng cũ. In giá trị nhỏ nhất và số vị trí đạt được.}{Vị trí tối ưu nằm tại trung vị của các tọa độ \(x\) và \(y\). Nếu \(N\) chẵn, tập tối ưu là một hình chữ nhật; trừ các vị trí đã có chuồng.}
\problemrow{Cow Sorting}{Cho một dãy, cần sắp tăng. Mỗi lần đổi chỗ hai số bất kỳ với chi phí bằng tổng hai số. Tối thiểu hóa tổng chi phí.}{Phân rã thành các chu trình. Với mỗi chu trình, so sánh hai cách: dùng phần tử nhỏ nhất trong chu trình để hoán đổi, hoặc đưa phần tử nhỏ nhất toàn dãy vào hỗ trợ.}
\problemrow{Lilypad Pond}{Ao \(M\times N\): \(1\) là đá, \(2\) là lá súng, \(3\) xuất phát, \(4\) đích. Chỉ nhảy theo bước mã và chỉ đáp xuống lá súng. Cần thêm ít lá nhất để đi từ đích đến xuất phát, đồng thời đếm số cách.}{Xem ô \(3,4\) như ô trống. Nếu từ một ô trống có thể qua vài ô đá rồi tới ô trống khác theo bước mã, nối cạnh giữa chúng; chạy đường ngắn nhất và đếm số đường.}
\problemrow{The Cow Lexicon}{Cho chuỗi dài \(L\) và \(N\) từ. Xóa ít ký tự nhất để chuỗi còn lại phân tách được thành các từ.}{\(F[i]\) là số ký tự xóa ít nhất để xử lý tiền tố \(i\). \(G[j+1][i]\) là số ký tự cần xóa trong đoạn để khớp một từ. Chuyển \(F[i]=\min_j(F[j]+G[j+1][i])\).}
\problemrow{Silver Cow Party}{Có \(N\) nông trại, \(M\) cạnh có hướng. Mỗi bò đi từ nông trại của mình tới \(P\), rồi quay về. Tìm thời gian lớn nhất trong các bò.}{Chạy Dijkstra từ \(P\) trên đồ thị gốc và trên đồ thị đảo cạnh; cộng hai khoảng cách cho từng đỉnh.}
\end{solutiontable}

\section{USACO 2007 January}

\begin{solutiontable}
\problemrow{Problem Solving}{Có \(P\) bài cần giải theo thứ tự. Bài \(i\) cần trả trước \(A_i\), sau khi giải còn trả \(B_i\). Mỗi tháng có \(M\) tiền. Hỏi ít nhất bao nhiêu tháng.}{DP \(F[i][j]\): đã giải \(i\) bài, tháng sau còn phải trả \(j\), số tháng ít nhất.}
\problemrow{Cow School}{Có \(N\) kỳ thi, kỳ \(i\) tổng điểm \(P_i\), đạt \(T_i\). Với mỗi \(D\), bỏ \(D\) tỷ lệ \(T_i/P_i\) nhỏ nhất rồi xét \(\sum T_i/\sum P_i\). Với mọi \(0\le D\le N\), hỏi có phương án khác tốt hơn không.}{Bài gốc chỉ ghi đề; cần so sánh các phương án bỏ bài theo tỷ lệ tổng, thường dùng biến đổi phân số và sắp xếp/kiểm tra tối ưu cho từng \(D\).}
\problemrow{Protecting the Flowers}{Bò \(i\) gây thiệt hại \(D_i\) mỗi đơn vị thời gian; đưa nó về chuồng mất \(T_i\), rồi quay lại mất \(T_i\). Tối thiểu hóa tổng thiệt hại.}{Tham lam sắp theo tỷ lệ \(T_i/D_i\) bằng so sánh chéo.}
\problemrow{Tallest Cow}{Có \(N\) bò, bò \(I\) cao nhất \(H\). \(R\) thông tin: bò \(A\) nhìn thấy bò \(B\), nghĩa là các bò giữa thấp hơn \(A\), và \(B\) không thấp hơn \(A\). Tìm chiều cao lớn nhất có thể của từng bò.}{Dùng mảng hiệu \(S\), ban đầu \(0\). Với thông tin \((A,B)\), chuẩn hóa \(A<B\), giảm \(S[A+1]\) và tăng \(S[B]\). Chiều cao bò \(i\) là \(H+\) tổng tiền tố tại \(i\).}
\problemrow{Balanced Lineup}{Có \(N\) bò cao \(H_i\), \(Q\) truy vấn hỏi hiệu giữa cao nhất và thấp nhất trong đoạn \([A,B]\).}{RMQ hoặc sparse table cho min và max.}
\end{solutiontable}

\section{USACO 2007 March}

\begin{solutiontable}
\problemrow{Gold Balanced Lineup}{Có \(N\) bò, \(K\) đặc trưng. Mỗi bò mã hóa tập đặc trưng bằng số. Đoạn cân bằng nếu mọi đặc trưng xuất hiện cùng số lần. Tìm đoạn cân bằng dài nhất.}{Đặt \(S[i][j]\) là số bò có đặc trưng \(j\) trong tiền tố \(i\). Đoạn \([A,B]\) cân bằng iff vector \(T[i][j]=S[i][j+1]-S[i][j]\) thỏa \(T[A-1]=T[B]\). Sắp hoặc hash các vector, giữ vị trí xuất hiện sớm nhất.}
\problemrow{Ranking the Cows}{Có \(N\) số và đã biết \(M\) quan hệ lớn nhỏ. Hỏi ít nhất cần so sánh thêm bao nhiêu cặp để xác định thứ tự.}{Tính bao đóng bắc cầu bằng BFS/Floyd. Với cặp chưa có quan hệ theo cả hai chiều, tăng đáp án.}
\problemrow{Face The Right Way}{Có \(N\) bò quay \(F\) hoặc \(B\). Mỗi thao tác đổi hướng \(K\) bò liên tiếp. Tìm \(K\) và số thao tác ít nhất để tất cả quay \(F\).}{Liệt kê \(K\). Với \(K\) cố định, quét trái sang phải; thao tác tại vị trí hiện tại bị quyết định bởi hướng hiện tại sau các lật trước đó.}
\problemrow{Cow Traffic}{DAG có cạnh \(i\to j\) với \(i<j\). Từ bất kỳ đỉnh bậc vào \(0\), đi tới \(N\). Tìm cạnh được đi qua nhiều nhất trong tất cả đường đi.}{\(F[i]\) là số đường từ nguồn bất kỳ tới \(i\); \(G[i]\) là số đường từ \(i\) tới \(N\). Cạnh \(i\to j\) được đi \(F[i]G[j]\) lần.}
\problemrow{Monthly Expense}{Chia \(N\) số thành \(M\) nhóm liên tiếp để tổng lớn nhất của một nhóm là nhỏ nhất.}{Nhị phân đáp án. Kiểm tra tham lam: đưa số vào nhóm hiện tại nếu không vượt ngưỡng, nếu không mở nhóm mới; khả thi nếu số nhóm không quá \(M\).}
\end{solutiontable}

\section{USACO 2007 November}

\begin{solutiontable}
\problemrow{Exploration}{Có \(N\) vị trí và \(T\) phút, xuất phát từ \(0\). Mỗi lần đi tới vị trí chưa thăm gần \(0\) nhất. Hỏi đi được bao nhiêu vị trí.}{Sắp theo \(|x|\), rồi đi lần lượt cho đến khi hết thời gian.}
\problemrow{Speed Reading}{Sách có \(N\) trang, \(K\) bò. Bò \(i\) đọc \(S_i\) trang/phút, đọc liên tục \(T_i\) phút rồi nghỉ \(R_i\) phút. Tính thời gian mỗi bò đọc xong.}{Mô phỏng chu kỳ đọc-nghỉ của từng bò.}
\problemrow{Avoid The Lakes}{Trong hình chữ nhật \(N\times M\), có \(K\) ô nước. Tìm kích thước thành phần nước liên thông lớn nhất.}{BFS hoặc DFS.}
\problemrow{Telephone Wire}{Có \(N\) số. Tổng chi phí gồm \(C\) nhân tổng \(|h_i-h_{i-1}|\); được tăng từng số, tăng \(X\) tốn \(X^2\). Tối thiểu hóa chi phí.}{DP \(F[i][j]\): chi phí nhỏ nhất khi số thứ \(i\) được nâng thành \(j\). Tách trường hợp \(j\ge h_i\) và \(j<h_i\), tối ưu xuống \(O(NH)\).}
\problemrow{Cow Relays}{Đồ thị vô hướng có \(T\) cạnh. Tìm đường ngắn nhất từ \(S\) đến \(E\) đi đúng \(N\) cạnh.}{Dùng lũy thừa ma trận trong đại số min-plus. \(D[k][i][j]\) là đường ngắn nhất từ \(i\) đến \(j\) dùng \(2^k\) cạnh; ghép bằng \(\min_w D[k][i][w]+D[k][w][j]\).}
\problemrow{Sunscreen}{Bò \(i\) dùng được SPF trong \([\min SPF_i,\max SPF_i]\). Có \(L\) chai SPF, mỗi chai dùng cho \(cover_i\) bò. Tối đa hóa số bò được bôi.}{Sắp chai theo SPF tăng dần; với mỗi chai, dùng cho các bò có thể dùng chai đó và có \(\max SPF\) nhỏ nhất.}
\problemrow{Cow Hurdles}{Đồ thị có hướng \(N\) đỉnh, \(M\) cạnh. Với mỗi truy vấn \(A_i,B_i\), tìm giá trị nhỏ nhất có thể của cạnh lớn nhất trên đường từ \(A_i\) đến \(B_i\).}{Biến thể Floyd: \(F[i][j]=\min(F[i][j],\max(F[i][k],F[k][j]))\).}
\problemrow{Milking Time}{Có \(N\) giờ và \(M\) khoảng vắt sữa \([A_i,B_i]\), nhận \(C_i\) sữa. Sau mỗi lần phải nghỉ \(R\) giờ. Tối đa hóa lượng sữa.}{Sắp khoảng theo \(A_i\). DP \(F[i]\) là lượng sữa tối đa đến khoảng \(i\), chuyển từ các \(j\) thỏa \(B_j+R\le A_i\).}
\problemrow{Best Cow Line}{Chuỗi dài \(N\). Mỗi lần lấy ký tự đầu hoặc cuối để tạo chuỗi kết quả nhỏ nhất theo thứ tự từ điển.}{Tham lam so sánh hai đầu; nếu bằng nhau thì tiếp tục so sâu vào trong.}
\end{solutiontable}

\section{USACO 2007 Open}

\begin{solutiontable}
\problemrow{Cheapest Palindrome}{Biến một chuỗi thành palindrome bằng chèn hoặc xóa ký tự với chi phí nhỏ nhất.}{Với ký tự \(c\), lấy \(w[c]\) là min giữa chi phí chèn và xóa. \(F[i][j]\) là chi phí nhỏ nhất để biến đoạn \(i..j\) thành palindrome: lấy min của xóa/chèn hai đầu, hoặc \(F[i+1][j-1]\) nếu hai đầu bằng nhau.}
\problemrow{Dining}{Có \(N\) bò, \(F\) loại thức ăn và \(D\) loại đồ uống. Mỗi bò thích một số thức ăn và đồ uống. Mỗi món chỉ cấp cho một bò; tối đa hóa số bò nhận được cả hai.}{Max flow: nguồn tới thức ăn, đồ uống tới đích; tách mỗi bò thành \(v_1\to v_2\); nối thức ăn bò thích tới \(v_1\), và \(v_2\) tới đồ uống bò thích, mọi dung lượng bằng \(1\).}
\problemrow{City Horizon}{Có các hình chữ nhật đáy trên trục \(x\) từ \(A_i\) tới \(B_i\), cao \(H_i\). Tính diện tích hợp.}{Sắp theo chiều cao giảm dần, dùng DSU hoặc cấu trúc đoạn để tô các đoạn chưa được phủ và cộng diện tích mới.}
\problemrow{Catch That Cow}{Từ số \(X\), mỗi bước được \(X-1\), \(X+1\), hoặc \(2X\). Hỏi ít bước nhất để biến \(N\) thành \(K\).}{BFS; bỏ trạng thái âm hoặc vượt quá \(2K\) theo ghi chú gốc.}
\problemrow{Fliptile}{Ma trận \(N\times M\) nhị phân. Bấm một ô sẽ lật ô đó và bốn ô kề. Tìm số lần bấm ít nhất để thành toàn \(0\), in phương án thứ tự từ điển nhỏ nhất.}{Liệt kê các ô được bấm ở hàng đầu; các hàng dưới bị quyết định bởi hàng trên. Cuối cùng kiểm tra hàng cuối toàn \(0\).}
\end{solutiontable}

\section{USACO 2008 February}

\begin{solutiontable}
\problemrow{Dining Cows}{Giống \emph{Eating Together}, nhưng giá trị chỉ là \(1\) hoặc \(2\), và cần biến dãy thành không giảm.}{DP đơn giản theo giá trị cuối cùng.}
\problemrow{Long Distance Racing}{Mỗi đoạn đường là lên dốc, phẳng hoặc xuống dốc. Thời gian tương ứng là \(U,F,D\). Trong thời gian cho trước, chạy được xa nhất rồi quay về xuất phát.}{Duyệt các đoạn theo thứ tự; nếu tổng thời gian đi và về cho đoạn tiếp theo còn đủ thì tiếp tục.}
\problemrow{Cow Multiplication}{Định nghĩa phép \(A*B\) là tổng tích của mọi cặp chữ số, một chữ số từ \(A\) và một chữ số từ \(B\). Tính \(A*B\).}{Cộng tổng chữ số của \(A\) và \(B\), đáp án là tích hai tổng đó.}
\problemrow{Making the Grade}{Cho dãy \(A_1,\ldots,A_N\). Tìm dãy \(B\) không tăng hoặc không giảm để tối thiểu hóa \(\sum |A_i-B_i|\).}{Luôn tồn tại nghiệm tối ưu mà mọi \(B_i\) là một giá trị xuất hiện trong \(A\). DP theo các giá trị đã sắp.}
\problemrow{Hotel}{Hỗ trợ hai thao tác: tìm vị trí đầu tiên có \(D\) phòng trống liên tiếp rồi lấp chúng; và làm trống \(D\) phòng liên tiếp bắt đầu từ \(X\).}{Cây đoạn lưu tiền tố trống, hậu tố trống và đoạn trống dài nhất, kèm lazy propagation.}
\problemrow{Game of Lines}{Trên mặt phẳng có \(N\) điểm; nối mọi cặp điểm tạo thành đường thẳng. Đếm số hướng không song song khác nhau.}{Tính hệ số góc chuẩn hóa cho mọi cặp, sắp xếp và loại trùng.}
\problemrow{Meteor Shower}{Có \(M\) sao băng; sao \(i\) rơi lúc \(T_i\) tại \((X_i,Y_i)\), từ thời điểm đó ô này và bốn ô kề không thể đi. Từ \((0,0)\), hỏi thời gian ít nhất để đến một ô an toàn vĩnh viễn.}{Tiền xử lý thời điểm nguy hiểm sớm nhất của từng ô, rồi BFS theo thời gian.}
\problemrow{Eating Together}{Có \(N\) số, mỗi số là \(1,2,3\). Đổi ít số nhất để dãy trở thành không giảm hoặc không tăng.}{DP theo vị trí và giá trị cuối cùng cho hai chiều, lấy min.}
\end{solutiontable}

\section{USACO 2008 January}

\begin{solutiontable}
\problemrow{Costume Party}{Có \(N\) bò, bò \(i\) có độ dài \(L_i\). Hai bò mặc chung một bộ nếu tổng độ dài không vượt \(S\). Hỏi ít nhất cần bao nhiêu bộ.}{Sắp \(L_i\) tăng dần. Mỗi lần xét bò lớn nhất còn lại; nếu ghép được với bò nhỏ nhất thì ghép, nếu không nó mặc riêng.}
\problemrow{Election Time}{Có \(N\) bò ứng viên, số phiếu kỳ vọng vòng một là \(A_i\), vòng hai là \(B_i\). Top \(K\) vòng một vào vòng hai; ai có \(B_i\) lớn nhất thắng. In chỉ số bò thắng.}{Mô phỏng: chọn \(K\) ứng viên có \(A_i\) lớn nhất, rồi chọn \(B_i\) lớn nhất trong nhóm đó.}
\problemrow{iCow}{Có \(N\) bài hát với giá trị \(R_i\). Mỗi vòng chọn bài có \(R_i\) lớn nhất, hòa thì chỉ số nhỏ nhất; chia đều giá trị đó cho \(N-1\) bài khác, phần dư cộng cho chỉ số nhỏ hơn, bài được chọn về \(0\).}{Mô phỏng đúng quy tắc phân phối.}
\problemrow{Artificial Lake}{Có \(N\) nền, mỗi nền có chiều cao và chiều rộng. Chọn nền thấp nhất rồi nhỏ nước từ đó; tính thời điểm mỗi nền bị ngập thêm \(1\)m.}{Mô phỏng quá trình nước dâng và tràn sang các nền lân cận theo thứ tự độ cao.}
\problemrow{Haybale Guessing}{Có \(N\) đống cỏ, số lượng đôi một khác nhau. Có \(M\) câu trả lời: min trên \([Q_l,Q_h]\) là \(A\). Tìm câu trả lời đầu tiên khiến các thông tin mâu thuẫn.}{Nhị phân số câu trả lời đúng ban đầu; với prefix giả định đúng, kiểm tra tính nhất quán bằng cấu trúc đoạn/DSU theo các ràng buộc min.}
\problemrow{Cell Phone Network}{Đặt ít trạm truyền thông nhất trên cây sao cho mọi đỉnh cách trạm gần nhất không quá \(1\).}{DP cây: \(F[i][0]\) là đỉnh \(i\) được con phủ, \(F[i][1]\) là đặt trạm tại \(i\), \(F[i][2]\) là được cha phủ.}
\problemrow{Cow Contest}{Có \(N\) bò trong cuộc thi với thứ hạng rõ ràng; \(M\) kết quả cho biết \(A\) hơn \(B\). Hỏi xác định được thứ hạng của bao nhiêu bò.}{Floyd/bao đóng bắc cầu quan hệ thắng-thua, rồi đếm bò có quan hệ xác định với mọi bò khác.}
\problemrow{Running}{Chạy \(N\) phút; phút \(i\) có thể chạy \(D_i\) mét. Mỗi phút chạy tăng mệt mỏi \(1\); nếu nghỉ phải nghỉ đến khi mệt mỏi về \(0\). Mệt mỏi không quá \(M\), cuối cùng phải về \(0\).}{DP \(F[i][j]\): sau phút \(i\), mệt mỏi \(j\), quãng đường lớn nhất. Chuyển bằng chạy hoặc nghỉ theo quy tắc.}
\problemrow{Telephone Lines}{Có \(N\) đỉnh, \(M\) cạnh. Được biến \(K\) cạnh trên đường đi thành chi phí \(0\). Tìm đường từ \(1\) đến \(N\) sao cho cạnh lớn nhất còn phải trả là nhỏ nhất.}{Nhị phân ngưỡng \(x\). Gán cạnh lớn hơn \(x\) chi phí \(1\), cạnh còn lại \(0\); tìm số cạnh lớn tối thiểu từ \(1\) tới \(N\), khả thi nếu không quá \(K\).}
\end{solutiontable}

\end{document}
